\documentclass[10pt,a4paper]{amsart}
\usepackage[margin=19mm]{geometry}
\usepackage{amsmath,amssymb,mathtools}
\usepackage[scr=rsfs,cal=euler]{mathalfa}
\usepackage{xfrac}
\usepackage[unicode,hidelinks,bookmarks=false]{hyperref}
\newcommand{\ie}{i.e.\ }
\newcommand{\cf}{cf.\ }
\newcommand{\resp}{respectively\ }
\newcommand{\Subsection}{\subsection}
\providecommand{\nobreakdash}{\nobreak\hbox{-}}
\setlength{\emergencystretch}{2em}
\multlinegap0pt
\newcommand\R{\mathbb R}
\usepackage{amssymb}
\newtheorem{proposition}{Proposition}
\title{Time warping with Hellinger elasticity}
\author{Yuly Billig}
\date{Source: arXiv:2603.08807v1 (CC BY 4.0)}
\begin{document}
\maketitle

\noindent\emph{Source and licence.} Time warping with Hellinger elasticity. Version arXiv:2603.08807v1 (CC BY 4.0), distributed by arXiv under the Creative Commons Attribution 4.0 International licence, http://creativecommons.org/licenses/by/4.0/, as recorded in the arXiv licence metadata for this exact version and shown on the abstract page https://arxiv.org/abs/2603.08807v1. Full licence notice: \texttt{licenses/arxiv-2603.08807-CC-BY-4.0.txt}.\par\smallskip

\noindent\textit{Editorial scope.} Counted unit: the article's proposition caseB (source0.tex lines 292--299). The article's complete original proof of it is delivered with it (source0.tex lines 300--320, one \texttt{proof} environment of 21 lines). No mathematics is written, repaired or re-derived: the statement and the proof are the article's own lines, in the article's own notation, and no retained line is substituted or re-flowed.  Retained uncounted as the source-local context the proof needs: lines 274--279.  Nothing was omitted from any retained span, so the package carries no omission record.  No retained line contains a citation, so the wrapper carries no bibliography and no bibliography entry.  No retained line contains a figure environment, an image-inclusion command or drawing code, so no image is delivered and the package declares no asset dependency.  Editorial formatting only, in addition: an AMS \texttt{amsart} preamble that restates the article's own declarations and macros used by the retained lines, namely the environments \texttt{proposition} on separate counters, together with the standard packages those lines need.  The retained environments print in this document as Proposition 1 in order of appearance, whereas the article prints the same items with the numbering of its own full document.  The complete native source of the article as distributed in the arXiv e-print for this exact version is bundled unchanged as \texttt{sources/196029/source0.tex}.
\begin{proposition}
\label{linear}
%Let $[a,b]$ be the intersection of segments $[\tau_i, \tau_{i+1}]$ and $[t_j, t_{j+1}]$. 
Let $[a,b] = [\tau_i, \tau_{i+1}] \cap [t_j, t_{j+1}]$. 
Then the optimal parametrization $\alpha$ is a linear function on $(a,b)$.
\end{proposition}

\begin{proposition}
\label{caseB}
Assume $\tau_i = t_j$ and $\tau_{i+1} = t_{j+k}$ for $k \geq 1$. Then for the optimal parametrization $\alpha$ with these constraints,
\begin{align*}
\int_{t_j}^{t_{j+k}}  C\big(f(\alpha(\tau)), g(\tau)\big) & \sqrt{\alpha^\prime(\tau)}  \,d\tau \\
= &\sqrt{(s_{i+1} - s_i) \sum_{\ell = 0}^{k-1} (t_{j+\ell+1} - t_{j+\ell}) C(f_i, g_{j+\ell})^2 }.
\end{align*}
\end{proposition}

\begin{proof}
By Proposition \ref{linear}, optimal parametrization $\alpha$ is piecewise linear. Denote by $k_\ell$ the slope of $\alpha(\tau)$ on interval $(t_{j+\ell}, t_{j+\ell+1})$. We want to maximize
\begin{equation*}
\sum_{\ell = 0}^{k-1} (t_{j+\ell+1} - t_{j+\ell}) C(f_i, g_{j+\ell}) \sqrt{k_\ell}
\end{equation*}
subject to the constraint
\begin{equation*}
\sum_{\ell = 0}^{k-1} (t_{j+\ell+1} - t_{j+\ell}) k_\ell = \int_{\tau_i}^{\tau_{i+1}} \alpha^\prime (\tau) \,d \tau = s_{i+1} - s_i.
\end{equation*}
For this, consider a positive definite scalar product on $\R^k$:
\begin{equation*}
({\bf x}, {\bf y}) = \sum_{\ell = 0}^{k-1} (t_{j+\ell+1} - t_{j+\ell}) \,x_\ell \,y_\ell,
\end{equation*}
and apply the Cauchy-Schwarz inequality for this product:
\begin{align*}
&\left( \sum_{\ell = 0}^{k-1}  (t_{j+\ell+1} - t_{j+\ell}) C(f_i, g_{j+\ell}) \sqrt{k_\ell} \right)^2 \\
&\quad\quad \leq \left( \sum_{\ell = 0}^{k-1} (t_{j+\ell+1} - t_{j+\ell}) k_\ell  \right) 
\left( \sum_{\ell = 0}^{k-1} (t_{j+\ell+1} - t_{j+\ell})  C(f_i, g_{j+\ell})^2 \right),
\end{align*}
and the equality holds when the slopes $k_\ell$ are proportional to $C(f_i, g_{j+\ell})^2$. The claim of the Proposition follows.
\end{proof}

\end{document}
